\documentclass{beamer}
\usepackage[utf8]{inputenc}
\usepackage{amsmath, amsfonts}
\usetheme{Madrid}
\title{Bilan de Stage : Modélisation et Pricing d'Options}
\author{Alexis VO}
\date{Juin 2025}

\begin{document}

\frame{\titlepage}

\begin{frame}{Objectifs du stage}
\begin{itemize}
    \item Étudier différents modèles de pricing d'options : CRR, Bachelier, Black-Scholes
    \item Implémenter un outil interactif en Python (Streamlit)
    \item Étudier les greeks et stratégies de couverture (hedging)
    \item Vulgariser les concepts dans des articles destinés à des étudiants sans pré-requis en finance
\end{itemize}
\end{frame}

\begin{frame}{Le modèle de Cox-Ross-Rubinstein (CRR)}
\begin{itemize}
    \item Modèle binomial discret pour valoriser des options
    \item Construction d'un arbre binaire avec facteurs $u$ et $d$
    \item Valeur de l’option obtenue par récurrence backward
    \item Convergence vers le modèle de Black-Scholes
\end{itemize}
\end{frame}

\begin{frame}{Implémentation technique (Python)}
\begin{itemize}
    \item Utilisation de \texttt{NumPy}, \texttt{matplotlib}, \texttt{Streamlit}
    \item Implémentation des modèles CRR, Black-Scholes, Bachelier
    \item Calcul des prix, greeks et stratégies de couverture
\end{itemize}
\end{frame}

\begin{frame}{Dashboard interactif}
\begin{itemize}
    \item Interface développée avec \texttt{Streamlit}
    \item Saisie des paramètres : $S$, $K$, $r$, $\sigma$, $T$
    \item Visualisation des résultats : prix, greeks, PnL
    \item Comparaison entre les modèles
    \item \textbf{Simulation en direct}
\end{itemize}
\end{frame}

\begin{frame}{Étude des Greeks et du Hedging}
\begin{itemize}
    \item Calculs analytiques ou numériques des greeks (delta, gamma, etc.)
    \item Implémentation du hedging dynamique (ex: couverture delta)
    \item Analyse de la performance selon les modèles
\end{itemize}
\end{frame}

\begin{frame}{Articles rédigés}
\begin{itemize}
    \item Article sur le modèle de CRR (fondements et application)
    \item Article sur le modèle de Bachelier
    \item Projet d’article sur le modèle de Black-Scholes
    \item Prix Fermat Junior
\end{itemize}
\end{frame}

\begin{frame}{Prix Fermat Junior}
\begin{itemize}
    \item Concours récompensant les meilleurs articles mathématiques d’étudiants (niveau Licence)
    \item Sujet envisagé : ?? sinon étude de la convergence CRR vers B-S
    \item Mise en valeur de ce stage
\end{itemize}
\end{frame}

\begin{frame}{Pistes pour la suite}
\begin{itemize}
    \item Article sur le modèle de Black-Scholes
    \item Implémentation de modèles plus complexes (Heston, SABR)
    \item Ajout de calibration, options exotiques, ou simulateur de PnL
\end{itemize}
\end{frame}

\begin{frame}{Conclusion}
\begin{itemize}
    \item Partie mathématique intéressante, implémentation enrichissante
    \item Outil interactif utile pour vulgariser les concepts
    \item Perspectives d'amélioration et de recherche
\end{itemize}
\end{frame}

\end{document}